\scanpage{53}
\mspage{6}
ou encore
\[
M'\times M'\xrightarrow{\lambda_G}\eta
\]
C'est obtenu ainsi : si $x,y\in M'(\unsure{\ell})$, on relève en
$\bar{x},\bar{y}\in\unsure{G|N}$, on prend $(\bar{x},\bar{y})=\bar{x}\bar{y}\bar{x}^{-1}\bar{y}^{-1}$, qui est
dans $N$, on réduit mod $N_0$, et c'est $\lambda_G(x,y)$.
(\deleted{\unsure{Comme}}\inserted{Si} $N'$ est central (ce qu'on suppose))
cette description montre que si $\unsure{\bar{x}}\ \unsure{\text{ou}|,}\ \unsure{\bar{y}}$
$\in d'(N')\subset M'$, alors \deleted{$\unsure{q}$} $\lambda_G(x,y)=0$, donc $\lambda_G$
s'interprète comme un hom
\[
\pi_0\otimes\pi_0\xrightarrow{\lambda_G}\eta=(N_{\mathrm{comm}})_M
\]
Notons aussi que $\lambda_G$ est alternée. Ceci
posé, \deleted{la condition} on a évidemment
\[
N/DG\simeq\eta/\lambda_G(\pi_0\otimes\pi_0)
\]
donc la condition cherchée est que
$\lambda_G$ \uline{est surjective}. Pour apprécier
cette condition, il faut voir comment
$\lambda_G$ varie avec $G$, on trouve (si $G_1$ \unsure{est à partir de|est déduit de}
$G$ défini par l'extension $E$ de $\pi_0$ par $\pi_1$)
\[
\lambda_{G_1}=\lambda_G+\beta c_E
\]
où $c_E:\pi_0\otimes\pi_0\longrightarrow\pi_1$ est l'accouplement
associé à l'ext. centrale de $\pi_0$ \inserted{(comm)} par $\pi_1$,
et où
\[
\beta:\pi_1\longrightarrow\eta=(N_{\mathrm{comm}})_M
\]
est le passage au quotient. \deleted{$\unsure{N_0}$}

Nous pouvons dévisser la condition
$\lambda_G$ \uline{surjective}, en introduisant encore
\mspage{52}
